\section{Turing Machine --- Build --- Step Clock}
%%% generated by the blueprint pipeline from TCSlib.Complexity.TuringMachine.Build.StepClock
\subsection*{Turing Machine --- Build --- Step Clock: Introduction}

Declarations of the Lean module \texttt{TCSlib.Complexity.TuringMachine.Build.StepClock}.

\subsection*{Turing Machine --- Build --- Step Clock: Results}

\begin{definition}[Configuration with an added clock tape]
\lean{Turing.clockCfg}
\label{Turing.clockCfg}
\leanok
\uses{Turing.Cfg}
Let $c$ be a configuration with $k$ work tapes over an alphabet $\Gamma$ and state type
$S$, on input $x$. Given a state value $q$ (possibly absent), a tape
$ct:\mathbb{Z}\to\mathrm{Option}(\Gamma)$, and a head position $cp\in\mathbb{Z}$, the
clock configuration has $k+1$ work tapes: it preserves the input-head position, the
original tapes and their head positions, and the output of $c$, while appending $ct$ as
the additional tape with head position $cp$ and using $q$ as its state.
\end{definition}

\begin{definition}[Word patch]
\lean{Turing.wordPatch}
\label{Turing.wordPatch}
\leanok
\uses{Turing.FinTM.bufferTape}
Given a tape $c_t\colon\mathbb{Z}\to\mathrm{Option}(\Gamma)$, an integer position $c_p$,
and a finite word $u$ over $\Gamma$, define a new tape that stores the symbols of $u$ in
the cells from $c_p$ through $c_p+|u|-1$, using the word’s positions relative to $c_p$,
and agrees with $c_t$ at every other cell.
\end{definition}

\begin{definition}[Stopped Turing-machine configurations]
\lean{Turing.Cfg.Stopped}
\label{Turing.Cfg.Stopped}
\leanok
\uses{Turing.Cfg}
For a natural number $k$, types $\Gamma$ and $S$, an input word $x\in\Gamma^*$, a
stopping predicate $\mathrm{stopAt}:S\to\{\mathrm{false},\mathrm{true}\}$, and a
configuration $c$ of type $\mathrm{Cfg}(k,\Gamma,S,x)$, say that $c$ is stopped when
either its state is absent or its state is $s\in S$ and $\mathrm{stopAt}(s)=\mathrm{true}$.
\end{definition}

\begin{definition}[Clock simulation states]
\lean{Turing.sc_State}
\label{Turing.sc_State}
\leanok
\uses{Turing.FlagPhase}
For a state set $S$, introduce a state type with four forms: a decrement state carrying
a body state and a flag phase; a body state carrying a body state; a finish state
carrying an optional body state; and a timeout state carrying an optional body state.
The flag phase is either the forward run phase, a rewind phase carrying a Boolean flag,
or a done phase carrying a Boolean flag.
\end{definition}

\begin{definition}[Decrement-phase exit transition]
\lean{Turing.sc_decExit}
\label{Turing.sc_decExit}
\leanok
\uses{Turing.FlagPhase, Turing.sc_State}
For a state $s$ and a flag phase, map the done phase with flag true to the body state
carrying $s$, and the done phase with flag false to the timeout state carrying $s$. Map
every other flag phase to the decrement state carrying $s$ and that phase.
\end{definition}

\begin{theorem}[Injectivity of the exit-state map]
\lean{Turing.sc_decExit_injective}
\label{Turing.sc_decExit_injective}
\leanok
\uses{Turing.sc_decExit}
\difficulty{4}
For every type $S$ and every fixed state $s\in S$, the map from counter phases to clock
states that sends the true completed phase to the body state at $s$, the false completed
phase to the timeout state carrying $s$, and every other phase to the decrement state
carrying $s$ and that phase is injective.
\end{theorem}

\begin{definition}[Clocked body action]
\lean{Turing.sc_bodyAct}
\label{Turing.sc_bodyAct}
\leanok
\uses{Turing.Action, Turing.sc_State}
Given an action for a machine with $k+1$ work tapes, form an action on the clock
simulation states that preserves its input-head movement, work-tape instructions, and
output symbol. Its successor state is the decrement state in the run phase carrying the
original successor state when that state exists, and the finish state carrying no body
state when the original action halts.
\end{definition}

\begin{definition}[Clocked machine construction]
\lean{Turing.sc_clockTM}
\label{Turing.sc_clockTM}
\leanok
\uses{Turing.Action.mapState, Turing.MultiTapeTM, Turing.MultiTapeTM.liftSymbols, Turing.decrementTM, Turing.embedEmitTM, Turing.lastEmb, Turing.sc_State, Turing.sc_bodyAct, Turing.sc_decExit}
Given an alphabet embedding that marks emitted symbols, a $k$-work-tape machine $M$, and
a Boolean stopping predicate on its states, construct a machine with $k+1$ work tapes
and the clock simulation states. Its initial state is the finish state carrying no body
state. In a decrement state, if the phase is run and the stopping predicate holds, it
enters the finish state carrying the current body state; otherwise it performs the
embedded decrement-machine transition and maps its successor state through the
decrement-phase exit map. In a body state it performs the embedded transition of $M$ and
maps the result to the clock simulation states. In finish and timeout states it stays in
the same state without moving heads, writing symbols, or emitting output.
\end{definition}

\begin{theorem}[Agreement of clock transitions]
\lean{Turing.sc_agree}
\label{Turing.sc_agree}
\leanok
\uses{Turing.Action.mapState, Turing.FlagPhase, Turing.MultiTapeTM, Turing.MultiTapeTM.liftSymbols, Turing.decrementTM, Turing.embedEmitTM, Turing.lastEmb, Turing.sc_bodyAct, Turing.sc_clockTM, Turing.sc_decExit}
\difficulty{4}
Let $\Gamma$ be an alphabet, let
$e:\{\mathrm{false},\mathrm{true}\}\hookrightarrow\Gamma$ be an embedding, and let $M$
be a deterministic Turing machine with $k$ work tapes over $\Gamma$. Let
$\mathrm{stopAt}$ assign a Boolean value to each state, and choose a state $s$ with
$\mathrm{stopAt}(s)=\mathrm{false}$. For any counter phase $f$ other than either
completed phase, any input symbol or blank, and any tuple of $k+1$ optional work-tape
symbols, the clock machine's transition from the state carrying $s$ and phase $f$ equals
the transition of the embedded Boolean decrement machine, with its successor state
relabeled by sending a completed true phase to the body state carrying $s$, a completed
false phase to the timeout state carrying $s$, and every other phase to the state
carrying $s$ and that phase.
\end{theorem}

\begin{theorem}[Layout of an embedded configuration]
\lean{Turing.sc_layout}
\label{Turing.sc_layout}
\leanok
\uses{Turing.Cfg, Turing.clockCfg, Turing.embedEmitCfg, Turing.embedEmitCfg_eq, Turing.sc_bodyAct}
\difficulty{4}
Let $c$ be a configuration of a $k$-work-tape machine on input $x$, and let $ct:\mathbb
Z\to\mathrm{Option}(\Gamma)$ and $cp\in\mathbb Z$. Embedding $c$ into the first $k$
tapes of a machine with one additional tape, assigning the additional tape content $ct$
and head position $cp$, and leaving the output prefix empty, gives exactly the clock
configuration with state $c$'s state and additional tape $(ct,cp)$.
\end{theorem}

\begin{theorem}[Layout of the final work tape]
\lean{Turing.sc_lastLayout}
\label{Turing.sc_lastLayout}
\leanok
\uses{Turing.Cfg, Turing.clockCfg, Turing.embedEmitCfg, Turing.embedEmitCfg_eq, Turing.lastEmb, Turing.lastEmb_apply, Turing.sc_bodyAct}
\difficulty{4}
Let $c$ be a configuration of a $k$-work-tape machine and $D$ a configuration with one
work tape, both on the same input. Let $ct:\mathbb Z\to\mathrm{Option}(\Gamma)$ and
$cp\in\mathbb Z$. If $D$ and $c$ have the same input-head position and output, then
embedding $D$'s sole work tape as the last tape among the $k+1$ tapes, using $c$'s tapes
and head positions for the first $k$ tapes, yields the clock configuration built from
$c$, with state and clock tape taken from $D$.
\end{theorem}

\begin{theorem}[Input-head position during decrement]
\lean{Turing.sc_dec_inputPos}
\label{Turing.sc_dec_inputPos}
\leanok
\uses{Turing.Action.apply, Turing.Cfg, Turing.FlagPhase, Turing.MultiTapeTM.runFrom, Turing.MultiTapeTM.runFrom_succ_eq_step', Turing.MultiTapeTM.step, Turing.decrementTM, Turing.decrementTM_tr_action, Turing.moveInputPos_zero}
\difficulty{4}
For every $k\in\mathbb{N}$, Boolean input word $x$, work-tape index
$i\in\mathrm{Fin}(k)$, configuration $d$ of the Boolean decrement machine on input $x$,
and time $t\in\mathbb{N}$, the input-head position after $t$ steps equals its initial
position in $d$.
\end{theorem}

\begin{theorem}[A word patch preserves an agreeing tape]
\lean{Turing.sc_wordPatch_self}
\label{Turing.sc_wordPatch_self}
\leanok
\uses{Turing.FinTM.bufferTape, Turing.wordPatch}
\difficulty{4}
Let $\Gamma$ be a set, let $c_t:\mathbb{Z}\to\operatorname{Option}(\Gamma)$ be a tape,
let $c_p\in\mathbb{Z}$, and let $u$ be a finite word over $\Gamma$. Suppose that for
every integer $p$ with $-1\le p\le |u|$, the tape content at $c_p+p$ equals the content
of $u$ at offset $p$. Patching the tape with $u$ at $c_p$ leaves the entire tape
unchanged.
\end{theorem}

\begin{theorem}[Repeated equal-length word patches]
\lean{Turing.sc_wordPatch_twice}
\label{Turing.sc_wordPatch_twice}
\leanok
\uses{Turing.wordPatch}
\difficulty{3}
Let $\Gamma$ be a set, let $c_t:\mathbb{Z}\to\operatorname{Option}(\Gamma)$ be a tape,
let $c_p\in\mathbb{Z}$, and let $u,v$ be finite words over $\Gamma$ of equal length.
Patching the tape first with $u$ and then with $v$, both at $c_p$, gives the same tape
as patching it once with $v$.
\end{theorem}

\begin{theorem}[Frame agreement for equal-length word patches]
\lean{Turing.sc_wordPatch_frame}
\label{Turing.sc_wordPatch_frame}
\leanok
\uses{Turing.FinTM.bufferTape, Turing.FinTM.bufferTape_left, Turing.wordPatch}
\difficulty{4}
Let $\Gamma$ be a set, let $c_t:\mathbb{Z}\to\operatorname{Option}(\Gamma)$ be a tape,
let $c_p\in\mathbb{Z}$, and let $u,v$ be finite words over $\Gamma$ of equal length.
Suppose that for every integer $p$ with $-1\le p\le |u|$, the tape content at $c_p+p$
equals the content of the word $u$ at offset $p$. Then replacing the tape on the
interval beginning at $c_p$ with $v$ gives, at every such offset $p$ from $-1$ through
$|v|$, exactly the content of $v$ at offset $p$.
\end{theorem}

\begin{theorem}[Width preservation under successful decrement]
\lean{Turing.sc_dec_success_width}
\label{Turing.sc_dec_success_width}
\leanok
\uses{Turing.counter_dec_spec, Turing.decFixed}
\difficulty{3}
For Boolean lists $w$ and $v$, if decrementing $w$ by the fixed-width decrement
operation returns $v$, then $v$ and $w$ have the same length.
\end{theorem}

\begin{theorem}[False words survive the prefix scan]
\lean{Turing.sc_takeWhile_replicate}
\label{Turing.sc_takeWhile_replicate}
\leanok
\difficulty{4}
For every $n\in\mathbb N$, the initial segment of $n$ false Boolean values has length
$n$ when scanned using the predicate that accepts false values.
\end{theorem}

\begin{theorem}[Cost of decrement underflow]
\lean{Turing.sc_dec_underflow_cost}
\label{Turing.sc_dec_underflow_cost}
\leanok
\uses{Turing.counter_dec_spec, Turing.decFixed, Turing.sc_takeWhile_replicate}
\difficulty{3}
For a Boolean list $w$, if the fixed-width decrement operation on $w$ is undefined, then
every entry of $w$ is false; equivalently, the initial segment of false entries in $w$
has length $|w|$.
\end{theorem}

\begin{theorem}[Width of the decremented word]
\lean{Turing.sc_dec_patch_width}
\label{Turing.sc_dec_patch_width}
\leanok
\uses{Turing.decFixed, Turing.sc_dec_success_width}
\difficulty{4}
For every Boolean word $w$, the list returned by decrementing $w$, with the all-true
word of length $|w|$ used when decrementing is undefined, has length $|w|$.
\end{theorem}

\begin{theorem}[Decoding preserves the Boolean counter tape]
\lean{Turing.sc_shadow_pre}
\label{Turing.sc_shadow_pre}
\leanok
\uses{Turing.Cfg, Turing.Cfg.decodeSymbols, Turing.Cfg.decodeSymbols_state, Turing.Cfg.decodeSymbols_workTapePos, Turing.Cfg.decodeSymbols_workTapes, Turing.FinTM.bufferTape, Turing.FinTM.bufferTape_map_decode, Turing.FlagPhase, Turing.decodeBool}
\difficulty{3}
Let $e$ inject the Boolean alphabet into $\Gamma$, let $x$ be an input over $\Gamma$,
let $ip$ be a valid input-head position, and let $out$ be an output word over $\Gamma$.
Let $w$ be a Boolean word, and suppose the tape $ct:\mathbb Z\to\mathrm{Option}(\Gamma)$
based at $cp\in\mathbb Z$ stores the encoded word $w$ at every offset $p$ with $-1\le
p\le |w|$. Form the configuration $D$ with input-head position $ip$, output $out$,
work-tape content $ct$, work-tape head position $cp$, and state $\mathrm{run}$. Decoding
$D$ leaves its state equal to $\mathrm{run}$, and its decoded work tape at each of those
offsets equals the buffer-tape content of $w$ there.
\end{theorem}

\begin{theorem}[Symbol lifting commutes with execution]
\lean{Turing.sc_lift_run}
\label{Turing.sc_lift_run}
\leanok
\uses{Turing.Cfg, Turing.Cfg.decodeSymbols, Turing.Cfg.liftOver, Turing.Cfg.liftOver_decodeSymbols, Turing.FlagPhase, Turing.MultiTapeTM.liftSymbols, Turing.MultiTapeTM.liftSymbols_runFrom_over, Turing.MultiTapeTM.runFrom, Turing.decrementTM, Turing.decrementTM_tr_action, Turing.decrementTM_tr_input}
\difficulty{3}
Let $e$ be an injection of the Boolean alphabet into an alphabet $\Gamma$, let $D$ be a
configuration of the one-tape machine over $\Gamma$ in the fixed input $x$, and let $t$
be a nonnegative integer. Running the symbol-lifted decrement machine from $D$ for $t$
steps gives the same configuration as first decoding the symbols of $D$, running the
Boolean decrement machine for $t$ steps, and then lifting that run over $D$ along $e$.
\end{theorem}

\begin{theorem}[Frame preservation under symbol lifting]
\lean{Turing.sc_lift_frame}
\label{Turing.sc_lift_frame}
\leanok
\uses{Turing.Cfg, Turing.Cfg.decodeSymbols, Turing.Cfg.decodeSymbols_inputPos, Turing.Cfg.liftOver_inputPos, Turing.Cfg.liftOver_output, Turing.Cfg.liftOver_workTapePos, Turing.FlagPhase, Turing.MultiTapeTM.liftSymbols, Turing.MultiTapeTM.output_prefix, Turing.MultiTapeTM.runFrom, Turing.decrementTM, Turing.sc_dec_inputPos, Turing.sc_lift_run}
\difficulty{4}
Let $e$ embed the Boolean alphabet into $\Gamma$, and let $D$ be a one-work-tape
configuration of the flag-phase machine on input $x$. If decoding the symbols of $D$ and
running the Boolean decrement machine for $T$ steps leaves its output empty, then at
every time $u\le T$ the lifted decrement machine on $\Gamma$ has the same input-head
position and output as $D$, and its work-tape head position agrees with that of the
decoded Boolean computation.
\end{theorem}

\begin{theorem}[Clocked Simulation of Boolean Decrement]
\lean{Turing.sc_debit}
\label{Turing.sc_debit}
\leanok
\uses{Turing.Action.mapState, Turing.Cfg, Turing.Cfg.decodeSymbols, Turing.Cfg.decodeSymbols_output, Turing.Cfg.decodeSymbols_workTapePos, Turing.Cfg.decodeSymbols_workTapes, Turing.Cfg.liftOver_state, Turing.Cfg.liftOver_workTapes, Turing.Cfg.mapState, Turing.FinTM.bufferTape, Turing.FinTM.bufferTape_map, Turing.FlagPhase, Turing.MultiTapeTM, Turing.MultiTapeTM.liftSymbols, Turing.MultiTapeTM.runFrom, Turing.MultiTapeTM.runFrom_mapState_of_agreeOn, Turing.clockCfg, Turing.decFixed, Turing.decodeBool, Turing.decodeBool_eq_some, Turing.decrementTM, Turing.decrementTM_run_succ_ofCfg, Turing.decrementTM_run_underflow_ofCfg, Turing.embedEmitCfg, Turing.embedEmitTM, Turing.embedEmitTM_runFrom, Turing.lastEmb, Turing.sc_dec_patch_width, Turing.sc_dec_underflow_cost, Turing.sc_lastLayout, Turing.sc_lift_frame, Turing.sc_lift_run, Turing.sc_shadow_pre, Turing.wordPatch}
\difficulty{7}
Let $k$ be a natural number, let $\Gamma,S,Q$ be types with decidable equality on
$\Gamma$, let $x$ be a finite word over $\Gamma$, let $e$ inject the Boolean alphabet
into $\Gamma$, let $H$ be a machine with $k+1$ work tapes, and let $\rho$ inject the
flagged phases into its state set $Q$. Suppose that for every phase other than a done
phase, every input symbol and every tuple of work-tape symbols, the transition of $H$ at
the corresponding state equals the transition obtained by embedding Boolean decrement on
the last tape, encoding symbols by $e$, and renaming phases by $\rho$. Let $c$ be a
configuration with $k$ work tapes on input $x$, let $w$ be a finite Boolean word, and
let $ct:\mathbb Z\to\mathrm{Option}(\Gamma)$ and $cp\in\mathbb Z$ satisfy that, for each
$p$ with $-1\le p\le |w|$, $ct(cp+p)$ equals the tape encoding of $w$ at offset $p$.
Starting with $c$ extended by this tape, its head at $cp$, and state $\rho(run)$, after
$2r+2$ steps, where $r$ is the number of initial false bits of $w$, the state is
$\rho(done((decFixed(w)).isSome))$, the original configuration components remain those
of $c$, and the added tape is patched with $e$ applied to the decremented word when
decrement is defined and otherwise to the all-true word of length $|w|$. Before that
time the machine is never in a done state. At every time up to and including $2r+2$, the
configuration consists of $c$ with its current state and added tape, the added head lies
in $[cp-1,cp+|w|]$, and its state is the image under $\rho$ of an optional flagged
phase.
\end{theorem}

\begin{theorem}[Stopping state triggers a finish step]
\lean{Turing.sc_stop_step}
\label{Turing.sc_stop_step}
\leanok
\uses{Turing.Cfg, Turing.MultiTapeTM, Turing.MultiTapeTM.step, Turing.clockCfg, Turing.moveInputPos, Turing.sc_clockTM}
\difficulty{1}
Let $k$ be a natural number, let $\Gamma$ and $S$ be types, let $e$ be an injection from
the Boolean alphabet into $\Gamma$, let $M$ be a Turing machine with $k$ work tapes,
alphabet $\Gamma$, and state set $S$, and let
$\mathrm{stopAt}:S\to\{\mathrm{false},\mathrm{true}\}$. For any finite input $x$ over
$\Gamma$, any configuration $c$ of $M$ on $x$, any state $s\in S$ with
$\mathrm{stopAt}(s)=\mathrm{true}$, any clock-tape contents $ct:\mathbb
Z\to\mathrm{Option}(\Gamma)$, and any clock-head position $cp\in\mathbb Z$, one step of
the clocked machine from the configuration with state $(s,\mathrm{run})$, body
configuration $c$, and clock data $ct,cp$ yields the configuration with state
$\mathrm{finish}(s)$ and the same body configuration and clock data.
\end{theorem}

\begin{theorem}[Clocked Body Step]
\lean{Turing.sc_body_step}
\label{Turing.sc_body_step}
\leanok
\uses{Turing.Action.apply, Turing.Cfg, Turing.Cfg.inputSymbol, Turing.Cfg.workTapeSymbols, Turing.MultiTapeTM, Turing.MultiTapeTM.runFrom, Turing.MultiTapeTM.step, Turing.clockCfg, Turing.embedEmitTM, Turing.embedEmitTM_runFrom, Turing.sc_bodyAct, Turing.sc_clockTM, Turing.sc_layout}
\difficulty{4}
Let $\Gamma$ be an alphabet, let
$e:\{\mathrm{false},\mathrm{true}\}\hookrightarrow\Gamma$ be an embedding, let $M$ be a
Turing machine with $k$ work tapes over $\Gamma$ and state set $S$, and let
$\mathrm{stopAt}:S\to\{\mathrm{false},\mathrm{true}\}$. For an input word $x$ over
$\Gamma$, a configuration $c$ of $M$ on $x$, and a state $s\in S$ such that $c$ is in
state $s$, let $ct:\mathbb{Z}\to\mathrm{Option}\,\Gamma$ be a tape content function and
$cp\in\mathbb{Z}$. One step of the clocked machine from the configuration with body
state $s$, underlying configuration $c$, and additional tape content $ct$ and head
position $cp$ yields the configuration whose underlying machine configuration is the
one-step successor of $c$, whose additional tape content and head position remain $ct$
and $cp$, and whose state is the decrement phase at the successor state if that
successor is a state, and the finish state if the successor configuration is halted.
\end{theorem}

\begin{definition}[Clock simulation exit states]
\lean{Turing.sc_exits}
\label{Turing.sc_exits}
\leanok
\uses{Turing.sc_State}
For the clock simulation state type, define the exit set to consist of precisely those
optional states that are some finish state or some timeout state, each carrying an
optional body state.
\end{definition}

\begin{definition}[Resume state]
\lean{Turing.sc_resume}
\label{Turing.sc_resume}
\leanok
\uses{Turing.sc_State}
For a body state $s$, map the optional state $\mathrm{some}(s)$ to the decrement state
carrying $s$ in the run phase, and map the absent state to the finish state carrying no
body state.
\end{definition}

\begin{theorem}[An observed computation segment]
\lean{Turing.sc_observedSegment}
\label{Turing.sc_observedSegment}
\leanok
\uses{Turing.Cfg, Turing.MultiTapeTM, Turing.MultiTapeTM.runFrom, Turing.MultiTapeTM.runFrom_succ_eq_step, Turing.MultiTapeTM.step, Turing.Segment, Turing.Segment.cons, Turing.Segment.zero}
\difficulty{4}
Let $N$ be a multi-tape Turing machine with state set $Q$, let
$F\subseteq\mathrm{Option}(Q)$ be a set of stopping states, let $P$ be a property of the
vector of work-tape head positions, and let $a$ be a configuration with fixed input $x$.
Suppose that at every time $u<t$ the state of the run from $a$ lies outside $F$, and
that at every time $u\le t$ the work-tape head positions satisfy $P$. Then the run from
$a$ for $t$ steps is a segment ending at the configuration reached at time $t$, avoiding
$F$ before time $t$ and satisfying $P$ through time $t$.
\end{theorem}

\begin{theorem}[Earlier states avoid the exit set]
\lean{Turing.sc_firstExit}
\label{Turing.sc_firstExit}
\leanok
\uses{Turing.Cfg, Turing.MultiTapeTM, Turing.MultiTapeTM.runFrom, Turing.Segment, Turing.Segment.take, Turing.Segment.time_eq}
\difficulty{3}
Let $N$ be a $k$-work-tape machine, let $F$ be a set of possible states (including the
halted state), and let $P$ be a property of the work-tape head positions. If a run of
$N$ from configuration $a$ reaches configuration $b$ after $t$ steps, avoids $F$ at
every time strictly before $t$, and satisfies $P$ at every time through $t$, then in
particular its state lies outside $F$ at every time strictly before $t$.
\end{theorem}

\begin{theorem}[Clocked decrement segment]
\lean{Turing.sc_clockDebit}
\label{Turing.sc_clockDebit}
\leanok
\uses{Turing.Cfg, Turing.FinTM.bufferTape, Turing.MultiTapeTM, Turing.MultiTapeTM.runFrom, Turing.Segment, Turing.clockCfg, Turing.decFixed, Turing.sc_agree, Turing.sc_clockTM, Turing.sc_debit, Turing.sc_decExit, Turing.sc_decExit_injective, Turing.sc_exits, Turing.sc_observedSegment, Turing.wordPatch}
\difficulty{5}
Let $e$ be an embedding of $\{0,1\}$ into an alphabet $\Gamma$, let $M$ be a
$k$-work-tape machine over $\Gamma$ with state set $S$, and let
$\mathrm{stopAt}:S\to\{\mathrm{false},\mathrm{true}\}$. Fix an input word $x$ over
$\Gamma$, a configuration $c$ of $M$ on $x$, a state $s\in S$ with
$\mathrm{stopAt}(s)=\mathrm{false}$, a Boolean word $w$, and a counter tape
$ct:\mathbb Z\to\mathrm{Option}(\Gamma)$ with head position $cp\in\mathbb Z$ such
that $ct(cp+p)$ contains the symbol of $w$ at offset $p$, mapped through $e$, for every
integer $p$ with $-1\le p\le |w|$ (and is blank where that word has no symbol). Form
the clocked machine $H$ from $M$, $e$, and $\mathrm{stopAt}$, and set
$t=2|w\mathbin{\mathrm{takeWhile}}(\mathrm{not})|+2$. After $t$ steps from the clock
configuration over $c$ with state recording the decrement phase for $s$ and counter tape
$ct$ at $cp$, the machine is in the clock configuration over $c$ with state determined
by whether the decrement of $w$ is defined, and counter tape obtained by patching $ct$
at $cp$ with the decremented word mapped through $e$ (using an all-true word of length
$|w|$ when the decrement is undefined). This run avoids the finish and timeout states at
every time strictly before $t$; at every time from $0$ through $t$, the counter head
lies in $[cp-1,cp+|w|]$, and the reached configuration is the clock configuration over
the unchanged $c$ with its current counter tape and head position.
\end{theorem}

\begin{theorem}[A configuration outside the stopping set is live]
\lean{Turing.sc_live}
\label{Turing.sc_live}
\leanok
\uses{Turing.Cfg, Turing.Cfg.Stopped}
\difficulty{4}
Let $c$ be a configuration of a $k$-tape Turing machine with state set $S$, and let
$\mathrm{stopAt}:S\to\{\mathrm{false},\mathrm{true}\}$ be a stopping predicate. If $c$
is not stopped—meaning its state is present and is not designated as a stopping
state—then there is a state $s\in S$ such that $c$ has state $s$ and
$\mathrm{stopAt}(s)=\mathrm{false}$.
\end{theorem}

\begin{theorem}[Clocked rounds and head control]
\lean{Turing.sc_rounds}
\label{Turing.sc_rounds}
\leanok
\uses{Turing.Cfg, Turing.Cfg.Stopped, Turing.FinTM.bufferTape, Turing.MultiTapeTM, Turing.MultiTapeTM.runFrom, Turing.MultiTapeTM.runFrom_add, Turing.MultiTapeTM.runFrom_succ_eq_step', Turing.MultiTapeTM.runFrom_zero, Turing.MultiTapeTM.step, Turing.Segment, Turing.Segment.one, Turing.Segment.trans, Turing.Segment.zero, Turing.clockCfg, Turing.counterOverhead, Turing.counterOverhead_succ, Turing.counterOverhead_zero, Turing.counterWord, Turing.counterWord_length, Turing.counterWord_zero, Turing.counter_word_decrement, Turing.sc_body_step, Turing.sc_clockDebit, Turing.sc_clockTM, Turing.sc_decExit, Turing.sc_exits, Turing.sc_live, Turing.sc_resume, Turing.sc_wordPatch_frame, Turing.sc_wordPatch_self, Turing.sc_wordPatch_twice, Turing.wordPatch}
\difficulty{6}
Let $e$ be an injection from $\{\mathrm{false},\mathrm{true}\}$ into an alphabet
$\Gamma$, let $M$ be a $k$-work-tape Turing machine with state set $S$, and let
$\mathrm{stopAt}:S\to\{\mathrm{false},\mathrm{true}\}$. Fix an input word $x$, a
configuration $c$ of $M$ on $x$, and a state $q\in S$ such that $c$ has state $q$; let
$w$ be a Boolean word, $ct:\mathbb Z\to\mathrm{Option}(\Gamma)$ a tape, and
$cp\in\mathbb Z$. Assume that for every integer $p$ with $-1\le p\le |w|$, the tape
cell $ct(cp+p)$ contains the symbol of $w$ at offset $p$ after encoding by $e$ (with
blank cells represented by absence). Let $n\in\mathbb N$ be at most the integer
represented by the binary digits of $w$, and suppose the run of $M$ from $c$ is not
stopped at any time $v<n$. Then, after $\mathrm{counterOverhead}(w,n)+n$ steps, the
clocked machine reaches the clock configuration whose body component is the
configuration of $M$ after $n$ steps, whose clock tape is patched with the encoded
counter word for round $n$, and whose clock head is at $cp$. The run segment up to that
time avoids the exit states before its endpoint and keeps the clock head in
$[cp-1,cp+|w|]$ throughout. Moreover, at every time $u\le
\mathrm{counterOverhead}(w,n)+n$, the clocked configuration consists of a clock state,
the body configuration of $M$ at some time $v\le n$, a tape, and a head position in
$[cp-1,cp+|w|]$.
\end{theorem}

\begin{theorem}[Delivery by the clocked machine]
\lean{Turing.sc_deliver}
\label{Turing.sc_deliver}
\leanok
\uses{Turing.Cfg, Turing.MultiTapeTM, Turing.MultiTapeTM.runFrom, Turing.Segment, Turing.clockCfg, Turing.sc_State, Turing.sc_clockTM, Turing.sc_exits, Turing.sc_firstExit, Turing.wordPatch}
\difficulty{4}
Let $e:\{\mathrm{false},\mathrm{true}\}\hookrightarrow\Gamma$ be an embedding, let $M$
be a machine with $k$ work tapes over $\Gamma$ and state set $S$, and let
$\mathrm{stopAt}:S\to\{\mathrm{false},\mathrm{true}\}$. Fix an input word $x$, a
configuration $c$ of $M$, a state $q\in S$, Boolean words $w,w'$, a tape content
$ct:\mathbb Z\to\mathrm{Option}(\Gamma)$, an integer head position $cp$, and natural
numbers $\tau,B,T$. Suppose $w'$ and $w$ have the same length, $T\le B$, the state
obtained by applying $\mathrm{exit}$ to the state of $M$ after $\tau$ steps belongs to
the finish-or-timeout exit set, and the clocked machine started from the configuration
with state $\mathrm{dec}(q,\mathrm{run})$, underlying configuration $c$, and clock tape
$(ct,cp)$ reaches after $T$ steps the configuration with state $\mathrm{exit}$ applied
to the state of $M$ after $\tau$ steps, underlying configuration $M$ after $\tau$ steps,
and clock tape obtained by patching $ct$ at $cp$ with the encoded word $e(w')$. Assume
this run avoids finish and timeout states before time $T$, its clock head remains
between $cp-1$ and $cp+|w|$ through time $T$, and at every time $u\le T$ its original
work-tape contents and head positions, input-head position, and output agree with those
of $M$ at some time $v\le\tau$. Then there are a time $T'\le B$ and a Boolean word $w''$
of length $|w|$ such that the clocked machine reaches the corresponding configuration
patched with $e(w'')$ at time $T'$, avoids finish and timeout states at every earlier
time, and at each time $u\le T'$ its original work tapes, input position, and output
agree with those of $M$ at some time $v\le\tau$, while its clock head remains between
$cp-1$ and $cp+|w|$.
\end{theorem}

\begin{theorem}[Bounded clock simulation]
\lean{Turing.sc_finish}
\label{Turing.sc_finish}
\leanok
\uses{Turing.Cfg, Turing.Cfg.Stopped, Turing.FinTM.bufferTape, Turing.MultiTapeTM, Turing.MultiTapeTM.runFrom, Turing.MultiTapeTM.runFrom_add, Turing.MultiTapeTM.step, Turing.Segment.one, Turing.Segment.trans, Turing.clockCfg, Turing.counterOverhead, Turing.counterOverhead_le_of_le, Turing.counterWord, Turing.counterWord_length, Turing.sc_clockTM, Turing.sc_deliver, Turing.sc_exits, Turing.sc_resume, Turing.sc_rounds, Turing.sc_stop_step, Turing.wordPatch}
\difficulty{5}
Let $\Gamma$ be an alphabet, let
$e:\{\mathrm{false},\mathrm{true}\}\hookrightarrow\Gamma$ be an embedding, and let $M$
be a deterministic machine with $k$ work tapes, state set $S$, and stopping predicate
$P:S\to\{\mathrm{false},\mathrm{true}\}$. Fix an input $x$, a configuration $c$ whose
state is $q$, a Boolean word $w$, a tape content $ct:\mathbb
Z\to\mathrm{Option}(\Gamma)$, and a position $cp\in\mathbb Z$. Suppose that for every
integer $p$ with $-1\le p\le |w|$, the content of this tape at $cp+p$ equals the content
prescribed there by the word $w$ encoded using $e$. For every natural number $\tau$ no
greater than the base-two value $\operatorname{Nat.ofDigits}(2,
w.map(\mathrm{Bool.toNat}))$, if $M$'s configuration after $\tau$ steps is stopped (its
state is absent or satisfies $P$), and every configuration at an earlier time is not
stopped, then the machine obtained by adjoining a counter tape and using it to clock the
simulation reaches, at some time $T\le 5\tau+2|w|+2$, a configuration whose simulated
machine configuration is exactly $M$'s configuration after $\tau$ steps, whose state is
the finish state carrying that configuration's state, and whose counter tape is patched
with the encoding under $e$ of some Boolean word $w'$ of length $|w|$. At every time
$u<T$ its state is neither a finish state nor a timeout state; at every time $u\le T$
its original work tapes, input-head position, and output agree with those of $M$ after
some number $v\le\tau$ of steps, and its counter-head position lies between $cp-1$ and
$cp+|w|$, inclusive.
\end{theorem}

\begin{theorem}[Bounded timeout of a clocked simulation]
\lean{Turing.sc_timeout}
\label{Turing.sc_timeout}
\leanok
\uses{Turing.Cfg, Turing.Cfg.Stopped, Turing.FinTM.bufferTape, Turing.MultiTapeTM, Turing.MultiTapeTM.runFrom, Turing.MultiTapeTM.runFrom_add, Turing.Segment.trans, Turing.clockCfg, Turing.counterOverhead, Turing.counterOverhead_le, Turing.counterOverhead_succ, Turing.counterWord, Turing.counterWord_length, Turing.counterWord_value, Turing.counter_dec_zero, Turing.decFixed, Turing.sc_clockDebit, Turing.sc_clockTM, Turing.sc_decExit, Turing.sc_deliver, Turing.sc_live, Turing.sc_resume, Turing.sc_rounds, Turing.sc_wordPatch_frame, Turing.sc_wordPatch_twice, Turing.wordPatch}
\difficulty{5}
Let $e:	ext{Bool}	o	ext{Γ}$ be an embedding, let $M$ be a $k$-work-tape machine over
$	ext{Γ}$ with state set $S$, and let $\mathrm{stopAt}:S\to\text{Bool}$. Let $x$ be a
finite input word and let $c$ be a configuration of $M$ on $x$ whose state is $q$. Let
$w$ be a Boolean word, let $ct:\mathbb Z\to\mathrm{Option}(\text{Γ})$ be a tape content
function, and let $cp\in\mathbb Z$ satisfy $ct(cp+p)$ equals the content at position $p$
of the tape storing the encoded word $w$ whenever $-1\le p\le |w|$. Write $N$ for the
integer represented by the binary digits of $w$, and suppose that for every natural
number $u\le N$, the configuration reached by running $M$ from $c$ for $u$ steps is not
stopped by $\mathrm{stopAt}$. Then there are a natural number $T\le 5N+2|w|+2$ and a
Boolean word $w'$ with $|w'|=|w|$ such that the clocked machine, started with body
configuration $c$, clock head at $cp$, and clock state carrying $q$ in its forward
counting phase, reaches after $T$ steps a timeout state carrying the state of $M$'s
configuration after $N$ steps, with the body tapes, input-head position, and output
equal to that configuration and the clock tape obtained by patching $ct$ with the
encoding of $w'$ at base position $cp$. At every time $u<T$, the clocked machine is in
neither its finish nor its timeout state. For every $u\le T$, there is a natural number
$v\le N$ such that its body tapes, their head positions, input-head position, and output
agree with those of $M$ after $v$ steps, and its clock-head position belongs to
$[cp-1,cp+|w|]$.
\end{theorem}

\begin{theorem}[Clocked simulation of a multi-tape machine]
\lean{Turing.exists_stepClockTMVia}
\label{Turing.exists_stepClockTMVia}
\leanok
\uses{Turing.Cfg, Turing.Cfg.Stopped, Turing.FinTM.bufferTape, Turing.MultiTapeTM, Turing.MultiTapeTM.runFrom, Turing.clockCfg, Turing.sc_State, Turing.sc_clockTM, Turing.sc_finish, Turing.sc_timeout, Turing.wordPatch}
\difficulty{3}
Let $\Gamma$ and $S$ be types, let $e:\{\mathrm{false},\mathrm{true}\}\hookrightarrow\Gamma$
be an embedding, let $M$ be a deterministic Turing machine with $k$ work tapes over
$\Gamma$ and state set $S$, and let $\mathrm{stopAt}:S\to\{\mathrm{false},\mathrm{true}\}$.
There exist a finite state set $Q$, a machine $H$ with $k+1$ work tapes over $\Gamma$,
an injective map $\mathrm{inj}:S\to Q$, and maps
$\mathrm{finish},\mathrm{timeout}:\mathrm{Option}(S)\to Q$ whose images are disjoint
and disjoint from the image of $\mathrm{inj}$. For every input $x\in\Gamma^*$,
configuration $c$ with state $q$, and Boolean word $w$ encoded on the clock tape at
offset $cp$ (with the specified one-cell margins), the construction has two behaviors.
If the body first reaches a configuration stopped by $\mathrm{stopAt}$ at time $\tau$
no later than the integer encoded by $w$, then within $5\tau+2|w|+2$ host steps it
reaches a finish state carrying the body's state, with the body configuration preserved
and the clock tape changed only on the encoded word's interval. Before that endpoint it
has not entered either exit state. At every time up to the endpoint, its body tapes,
head positions, input head, and output agree with the body machine at some time at most
$\tau$, and the clock head stays in $[cp-1,cp+|w|]$. If the body never reaches a
stopped configuration by the value encoded by $w$, the analogous conclusion holds with
a timeout state and the body configuration at that encoded time.
\end{theorem}

\begin{theorem}[Finite-time clock simulation]
\lean{Turing.exists_stepClockTM}
\label{Turing.exists_stepClockTM}
\leanok
\uses{Turing.Cfg, Turing.FinTM.bufferTape, Turing.MultiTapeTM, Turing.MultiTapeTM.runFrom, Turing.clockCfg, Turing.exists_stepClockTMVia}
\difficulty{4}
Let $M$ be a deterministic Turing machine with $k$ work tapes, Boolean alphabet, and
finite state set $S$, and let $P:S\to\{\mathrm{false},\mathrm{true}\}$ be a state
predicate. There exist a finite-state machine $H$ with $k+1$ work tapes, an injective
map $\iota:S\to Q$, and maps $F,U:\mathrm{Option}(S)\to Q$ such that the following
holds. For every Boolean input $x$, every initial configuration $c$ of $M$ whose state
is $q$, every Boolean word $w$, and every clock-tape offset $a\in\mathbb Z$, suppose the
extra tape content $C$ satisfies $C(a+p)$ equals the content of the word $w$ at position
$p$ for each integer $p$ with $-1\le p\le |w|$. Set $d$ to the natural number
represented by the base-two digits of $w$. If, at some time $\tau\le d$, the run of $M$
from $c$ is halted or its state satisfies $P$, and every earlier configuration in that
run is neither halted nor in a state satisfying $P$, then for some $T\le 5\tau+2|w|+2$,
machine $H$, started from the configuration encoding $c$ with state $\iota(q)$ and extra
tape $C$ at offset $a$, reaches the configuration of $M$ at time $\tau$ with state $F$
applied to that configuration's state; before time $T$, its state is never in the image
of $F$ or $U$. If no configuration of $M$ at any time $u\le d$ is halted or has a state
satisfying $P$, then for some $T\le 5d+2|w|+2$, the same run of $H$ reaches the
configuration of $M$ at time $d$ with state $U$ applied to that configuration's state;
before time $T$, its state is never in the image of $F$ or $U$.
\end{theorem}
